\documentclass[10pt,a4paper]{amsart}
\usepackage[margin=19mm]{geometry}
\usepackage{amsmath,amssymb,mathtools}
\usepackage[scr=rsfs,cal=euler]{mathalfa}
\usepackage{xfrac}
\usepackage[unicode,hidelinks,bookmarks=false]{hyperref}
\newcommand{\ie}{i.e.\ }
\newcommand{\cf}{cf.\ }
\newcommand{\resp}{respectively\ }
\newcommand{\Subsection}{\subsection}
\providecommand{\nobreakdash}{\nobreak\hbox{-}}
\setlength{\emergencystretch}{2em}
\multlinegap0pt
\usepackage{bm}
\usepackage{bbm}
\usepackage{dsfont}
\usepackage{xspace}
\usepackage{cancel}
\usepackage{mathtools}
\usepackage{amsthm}
\makeatletter
\@ifundefined{theorem}{\newtheorem{theorem}{Theorem}}{}
\@ifundefined{corollary}{\newtheorem{corollary}[theorem]{Corollary}}{}
\@ifundefined{proposition}{\newtheorem{proposition}[theorem]{Proposition}}{}
\makeatother
\newcommand{\R}{\mathbb{R}}
\title[An Algebraic Representation Theorem for Linear GENEOs in Geo]{An Algebraic Representation Theorem for Linear GENEOs in Geometric Machine Learning}
\author{Francesco Conti, Patrizio Frosini, Nicola Quercioli}
\date{Source: arXiv:2601.03910v2 (CC BY 4.0)}
\begin{document}
\maketitle

\noindent\emph{Source and licence.} An Algebraic Representation Theorem for Linear GENEOs in Geometric Machine Learning. Version arXiv:2601.03910v2 (CC BY 4.0), distributed under the Creative Commons Attribution 4.0 International licence, as recorded in the arXiv licence metadata for this exact version and shown on the abstract page https://arxiv.org/abs/2601.03910v2. Full rights notice: licenses/arxiv-2601.03910-CC-BY-4.0.txt.\par\smallskip

\noindent\textit{Editorial scope.} Counted unit: the Corollary whose source label is \texttt{cor:convexity} in the article (source0.tex lines 1052--1058, source label \texttt{cor:convexity}), delivered together with the article's own complete original proof of it (source0.tex lines 1060--1153, one \texttt{proof} environment of 94 lines). No mathematics is written, repaired or re-derived: statement and proof are the article's own text line for line. Required context retained uncounted: maint (lines 310--313); teo:linear (lines 781--784); prop:convexity (lines 1034--1040). Every definition, display and label of those lines is retained unchanged. Omitted from this package and not redistributed: 1--309, 314--780, 785--1033, 1041--1051, 1059--1059, 1154--1565. Nothing inside a retained excerpt is omitted or altered: every retained body is delivered exactly as the article's lines are written, so no omission receipt of any kind accompanies this package. Citations: the retained excerpts carry no citation command at all, so no bibliography entry of the article is transcribed and no reference is redistributed. Reference adaptation: none was needed and none was made; every reference inside the delivered unit resolves inside it, so no unresolved reference or citation is produced. Printed slips kept literal: no printed word, symbol, hypothesis or number of the counted statement or of its proof was altered. Layout and class: the article's own class is \texttt{article} with packages including geometry, amsthm, graphicx, amsmath, amssymb, amscd, mathtools, dsfont, xcolor, todonotes, ulem, hyperref, cancel; the wrapper is the lane's pinned \texttt{amsart} editorial wrapper at 10pt on A4, which restates the theorem-like environments the retained text needs and adds the article's own macro definitions that the retained text needs (\texttt{\textbackslash R}), with extra wrapper packages (bm, bbm, dsfont, xspace, cancel, mathtools). The delivered numbering is the unit's own. Assets: no retained span contains a figure environment, a graphics-inclusion command, a PDF-page inclusion, a table or a TikZ picture; no image file is copied into this package, the package declares no asset dependency and no image file is redistributed with it. Licence: version arXiv:2601.03910v2 of this article is distributed under the Creative Commons Attribution 4.0 International licence, as recorded in the arXiv licence metadata for this exact version and shown on the abstract page https://arxiv.org/abs/2601.03910v2. A full rights notice is delivered at licenses/arxiv-2601.03910-CC-BY-4.0.txt. Characters: every retained excerpt is pure ASCII, so the delivered engine, XeLaTeX with its default Latin Modern text font, reports no missing character.
\begin{theorem}
    \label{maint} 
    Assume that $G \subseteq \mathrm{Aut}(X), K \subseteq \mathrm{Aut}(Y)$, and $T \colon G \to K$ is a group homomorphism. Moreover, suppose that $T(G)$ transitively acts on the finite set $Y$, and $F$ is a map from $\mathbb{R}^X$ to $\mathbb{R}^Y$. The map $(F,T)$ is a linear GEO from $(\mathbb{R}^X, G)$ to $(\mathbb{R}^Y, K)$ if and only if a $T$-permutant measure $\mu$ exists such that $F(\varphi) = \sum_{h \in X^Y} \varphi h\  \mu(h)$ for every $\varphi \in \mathbb{R}^X$.
\end{theorem}

\begin{theorem}
    \label{teo:linear}
    Assume that $G \subseteq \mathrm{Aut}(X), K \subseteq \mathrm{Aut}(Y)$, $T \colon G \to K$ is a group homomorphism, $T(G) \subseteq K$ transitively acts on the finite set $Y$ and $F$ is a map from $\mathbb{R}^X$ to $\mathbb{R}^Y$. The map $(F,T)$ is a linear GENEO from $(\mathbb{R}^X, G)$ to $(\mathbb{R}^Y, K)$ if and only if a generalized $T$-permutant measure $\mu$ exists such that $F(\varphi) = \sum_{h \in X^Y} \varphi h \ \mu(h)$ for every $\varphi \in \mathbb{R}^X$ and $\sum_{h \in X^Y} \left| \mu(h) \right| \le 1$.
\end{theorem}

\begin{proposition}
\label{prop:convexity}
The signed measures $\nu_1, \ldots, \nu_k$ form a basis for $\Gamma$.
% The $T$-permutant measures on $X^Y$ are exactly the signed measures that can be expressed as linear combinations of the signed measures $\nu_1, \ldots, \nu_k$.
    %If $\mu$ is a $T$-permutant measure, then
    %$\mu = \sum_{i=1}^k \mu(\mathcal{O}_i) \nu_i$.
\end{proposition}

\begin{corollary}
\label{cor:convexity}
The linear GEOs from $(\R^X,G)$ to $(\R^Y,K)$ are exactly the pairs of maps  
$(\sum_{i=1}^k a_i F_{\nu_i},T)$, with $a_1,\ldots,a_k$ arbitrary real numbers. The linear GENEOs from $(\R^X,G)$ to $(\R^Y,K)$ are exactly the pairs of maps $(\sum_{i=1}^k a_i F_{\nu_i},T)$, with  $\sum_{i=1}^k |a_i||\mathcal{O}_i|\le 1$.
    %If $\mu$ is a $T$-permutant measure, then
    %$\mu = \sum_{i=1}^k \mu(\mathcal{O}_i) \nu_i$.
\end{corollary}

\begin{proof}
Theorem~\ref{maint}  guarantees that the linear GEOs from
$(\mathbb{R}^X, G)$ to $(\mathbb{R}^Y, K)$ are precisely the pairs
$(F_\mu, T)$, where $\mu$ is a generalized $T$-permutant measure on $X^Y$.
The first statement of Corollary~\ref{cor:convexity} follows immediately from
this result and Proposition~\ref{prop:convexity}, by observing that if
$\mu = \sum_{i=1}^k a_i \ \nu_i$, then
$F_\mu = \sum_{i=1}^k a_i F_{\nu_i}$.
%since the orbits $\mathcal{O}_i$ are disjoint.
Moreover, Theorem~\ref{teo:linear} guarantees that the linear GENEOs from
$(\mathbb{R}^X, G)$ to $(\mathbb{R}^Y, K)$ are exactly the pairs
$(F_\mu, T)$ for which $\mu$ is a $T$-permutant measure on $X^Y$
satisfying $\sum_{h \in X^Y} |\mu(h)| \le 1$. The second statement of Corollary~\ref{cor:convexity}
follows from Theorem~\ref{teo:linear} by observing that, if $\mu=\sum_{i=1}^k a_i \ \nu_i$, then 

\[
\begin{aligned}
\sum_{h\in X^Y} |\mu(h)|
&=
\sum_{h\in X^Y} \left|\sum_{j=1}^k a_j \ \nu_j(h)\right|\\
&=\sum_{i=1}^k \sum_{h\in \mathcal{O}_i} \left|\sum_{j=1}^k a_j \ \nu_j(h)\right|\\
&=\sum_{i=1}^k \sum_{h\in \mathcal{O}_i} \left|a_i \right|\\
&=\sum_{i=1}^k |a_i| |\mathcal{O}_i|.
\end{aligned}
\]
%\hrule
%Let us start from the statement about GEOs.
% On the one hand, it is trivial to observe that
% $\left(\sum_{i=1}^k a_i F_{\nu_i}, T\right)$
% is a GEO from $(\mathbb{R}^X, G)$ to $(\mathbb{R}^Y, K)$.
% If $\sum_{i=1}^k |a_i| |\mathcal{O}_i|\le 1$, then
% $\left(\sum_{i=1}^k a_i F_{\nu_i}, T\right)$
% is a GENEO, since
% \[
% \begin{aligned}
% \left\|
% \sum_{i=1}^k a_i F_{\nu_i}(\p_1)
% - \sum_{i=1}^k a_i F_{\nu_i}(\p_2)
% \right\|_\infty
% &=
% \left\|
% \sum_{i=1}^k a_i F_{\nu_i}(\p_1 - \p_2)
% \right\|_\infty \\
% &\le
% \sum_{i=1}^k |a_i|
% \left\|F_{\nu_i}(\p_1 - \p_2)\right\|_\infty \\
% %&= \sum_{i=1}^k |a_i|
% %\left\|(\p_1 - \p_2)h \nu_i(h)\right\|_\infty \\
% &\le
% \sum_{i=1}^k |a_i| |\mathcal{O}_i|\,
% \left\|\p_1 - \p_2\right\|_\infty .
% \end{aligned}
% \]

% On the other hand, if $(F, T)$ is a GEO from $(\mathbb{R}^X, G)$ to $(\mathbb{R}^Y, K)$,
% then Theorem~\ref{teo:linear} implies the existence of a generalized $T$-permutant measure $\mu$
% such that
% $F(\varphi) = \sum_{h \in X^Y} \varphi h \, \mu(h)$
% for every $\varphi \in \mathbb{R}^X$.
% From Proposition~\ref{prop:convexity}, there exists $(a_1, \ldots, a_k) \in \mathbb{R}^k$
% such that
% $\mu = \sum_{i=1}^k a_i \nu_i$,
% and hence
% \begin{align*}
% F(\varphi)
% &= \sum_{h \in X^Y} \varphi h \sum_{i=1}^k a_i \nu_i(h) \\
% &= \sum_{i=1}^k a_i \sum_{h \in X^Y} \varphi h \nu_i(h)
% = \sum_{i=1}^k a_i F_{\nu_i}(\varphi).
% \end{align*}
% If $(F,T)$ is a GENEO, then $F$ is non-expansive and hence, in particular, 
% $\|\sum_{i=1}^k a_i F_{\nu_i}(\varphi)\|_\infty\le \|\p\|_\infty$,
% i.e.,
% $\|\sum_{i=1}^k a_i \sum_{h\in \mathcal{O}_i} \varphi h\|_\infty\le \|\p\|_\infty$ for every $\p\in \R^X$.
% If we set $\p={\mathds{1}}_{x_j}$,
% we obtain the inequality
% $\|\sum_{i=1}^k a_i \sum_{h\in \mathcal{O}_i} {\mathds{1}}_{x_j} h\|_\infty\le \|{\mathds{1}}_{x_j}\|_\infty=1$.






% Moreover,
% $\|F(\varphi)\|_\infty
% \le \sum_{i=1}^k |a_i| \, \|F_{\nu_i}(\varphi)\|_\infty
% \le \sum_{i=1}^k |a_i| \, |\mathcal{O}_i| \, \|\varphi\|_\infty$
% for every $\varphi \in \mathbb{R}^X$.
% Since $F$ is linear, if
% $\sum_{i=1}^k |a_i| \, |\mathcal{O}_i| \le 1$,
% then
% $\|F(\varphi_1 - \varphi_2)\|_\infty
% \le \|\varphi_1 - \varphi_2\|_\infty$
% for every $\varphi_1, \varphi_2 \in \mathbb{R}^X$.
\end{proof}

\end{document}
